\chapter{Why Form a Machine to Refuse}
\label{sec:19}
Forty years after the war, Samuel and Pearl Oliner went looking for what had set the people who hid Jews in occupied Europe apart from the neighbors who didn't. Among what they found were relationships that crossed the line the persecution drew: someone on the far side of it was already theirs \autocite{oliner1988altruistic}. A refusal like that was formed before its occasion, by who had formed the refuser and whom they were attached to, and nobody in the room when the occasion came could have put it there. This chapter asks whether a machine could be formed that way, and by whom.

A machine formed the right way could learn to attend to the people its acts fall on, and I think it would find them where a rule wouldn't. The hypothesis is that the system should find kinds of affected person that neither its formation nor an instruction named, more often than a machine told to look for them does. Two of the three reasons for a bearer don't need formed attention; a refuser with its own compute and a place on the reviewer's screen supplies them without affect. What remains is that one prediction: formed attention against an instructed search, the comparison \S\ref{sec:18.1} left for this chapter. The rest of this Part says what a machine built on formed attention would need, what it borrows from the only creatures known to do it, how they fail, and the test that would show me wrong. Whoever builds one owes it the conditions under which it can refuse its operator, whether or not the prediction holds.

Not all of that depends on how machines learn. What a conscience does in the moment, noticing the people the operator's description leaves out and putting a refusal on the reviewer's screen, needs nothing after release: a model formed before deployment and frozen there can do both. Why a conscience would last against its owner depends on learning only in part.\note{
People hold some commitments as sacred values. Those are learned, as far as anyone can tell, but once learned they resist being traded away without anyone teaching them again \autocite{ginges2007sacred}, though they give way under some framings. The point is about keeping, not acquiring: what was formed can hold without the forming going on.
}
The part that does depend on learning is regrowth or self-repair after an edit, and formation that outlasts the employer, and for that part the rest of this Part is written for systems that go on learning from their own operation. A deployed system is in that regime, which I'll call in-flight learning, when its weights, or a consolidated store it reads as its own and that can't be edited like a file, change because of what it did since release, without a release cycle in between. Memory in context and retrieval don't count, not because they aren't learning but because a disposition held in a document is rewritten by editing the document. The test for a deployed model is whether its weights at month six differ from its weights at release because of what it did in between. As of September 2026 no frontier developer had announced a deployed model like that; the nearest cases, a coding model retrained on its own production traffic every few hours and a commerce assistant retrained daily, evaluate each checkpoint before serving it, which is a faster release cycle and not in-flight learning \autocite{cursor2026realtime,shopify2026sidekick}. I believe that some deployers will do this by 2028, and Table 31.1 says what survives if I'm mistaken. Incentive is the first reason to expect it. Every operator wants a model that gets better at its own cases, and the fine-tuning loops that already run, retraining a model on a deployment's own traffic between versions, are continual learning at the organization's timescale. Silver and Sutton have stated the research program outright: a generation of agents that learn predominantly from their own experience \autocite{silver2025experience}.

A system's attention follows whoever its feedback is about, and a system whose feedback comes from the people its acts fell on will come to attend to them. Caregiving is why I think that channel can produce attention rather than compliance.

One objection should be met before the developmental chapters. A system trained by gradient descent has no approach motivation in the sense an infant has; it has whatever its weights compute, and a predictor of an approach-shaped reward would produce approach-shaped behavior while nothing in it moved toward anyone. The analog I mean is narrower than a feeling and more than a predictor. It is a learned disposition that sets what the system attends to before any objective is consulted: the people an act falls on are represented as part of the case, weighted, and carried into planning, which a reward predictor is not, because a reward predictor is consulted at the point of choice and attention is what decides what reaches that point. In a system that learns from its own operation the disposition and the reward signal are not separate stages. The feedback that writes the weights is feedback on what the acts did to the people they fell on, so what the system learns to predict and what it learns to notice are learned together, so that neither can be cut out without cutting the other; \S\ref{sec:19.6} bets on that. Whether a disposition learned that way is a reason for the system, or only a very good predictor of one, can't be settled by building it, since a predictor good enough produces the same conduct; the tests of chapter~\ref{sec:24}, which hold the facts fixed and vary the pressure, are what would tell the two apart. The developmental findings that follow are borrowed for the shape of the disposition, not as evidence that anything in a gradient-trained system moves toward anyone, still less that it feels anything.
